\section{Analysis}
\label{sec:analysis}

\subsection{Failure Mode Taxonomy}

Based on the pilot evaluation, we define six failure modes specific to constructive algebraic geometry reasoning.
While the high overall accuracy (78\%) and the dominance of timeouts among failures limit our quantitative observations, the taxonomy provides a systematic framework for characterizing model weaknesses.

\paragraph{F1: Fabricated objects.}
The model produces a mathematical object that does not exist or is not well-defined.
\emph{Definition:} CAS verification reveals that the stated object does not have the claimed algebraic structure (e.g., a ``smooth curve'' that is singular, a ``coherent sheaf'' with incompatible transition functions).

\paragraph{F2: Property confusion.}
The model constructs a valid object but confuses the required properties---either verifying the wrong property, or failing to distinguish between subtly different conditions.
\emph{Pilot example (\ecagid{0343}):} In a Level~3 problem about a $\mathbb{Z}/2\mathbb{Z}$-action on $\mathbb{P}^3$, the model substituted a different group action from the one specified, then correctly verified properties for this substituted action.
The algebraic computations were sound, but the mismatch with the intended construction constitutes a property confusion error.

\paragraph{F3: Computational errors.}
The model follows a correct strategy but makes arithmetic or algebraic errors in computation.
This class of errors is particularly insidious because the reasoning structure is correct, and only explicit CAS computation reveals the mistake.

\paragraph{F4: Trivial or degenerate answers.}
The model produces a technically correct but trivially degenerate answer that demonstrates no mathematical understanding.
\emph{Definition:} Answers that use identity morphisms, zero sheaves, or other universal constructions that trivially satisfy conditions without illustrating the mathematical concept.

\paragraph{F5: Hallucinated theorems.}
The model invokes a theorem that does not exist or misattributes a result.
This failure mode is expected to be concentrated at higher difficulty levels, where the model is more likely to encounter knowledge gaps that it fills with plausible-sounding fabrications.

\paragraph{F6: Incomplete verification.}
The model constructs a plausible object and provides some justification but fails to verify all required properties.
\emph{Pilot example (\ecagid{0095}):} On a Level~3 problem about blow-ups, the model provided a textbook-perfect construction of the blow-up of $\mathbb{A}^2$ at the origin, but the problem asked for a more specific construction that the model did not fully address, resulting in partial credit.

\paragraph{Observed failure mode distribution.}
Among the 82 problems that received responses, only 4 were scored below 1.0.
Of these, 3 received F6 (incomplete verification) and 1 received F2 (property confusion).
The remaining 18 failures were timeouts, for which no failure mode can be assigned since the model produced no response.
This sparsity limits failure mode analysis and suggests that future evaluations with larger samples and multiple models will be needed to fully populate the taxonomy.

\subsection{Difficulty Calibration}
\label{sec:difficulty-calibration}

The pilot results show a clear monotonic relationship between assigned difficulty and model accuracy:

\begin{center}
\begin{tabular}{lccc}
\toprule
& Level 1 & Level 2 & Level 3 \\
\midrule
$n$ & 22 & 67 & 11 \\
Accuracy & 0.95 & 0.81 & 0.27 \\
Avg.\ score & 0.95 & 0.81 & 0.41 \\
Timeout rate & 4.5\% & 17.9\% & 45.5\% \\
\bottomrule
\end{tabular}
\end{center}

This validates the difficulty calibration: problems rated harder by a human expert are indeed substantially harder for the model.
The steepest drop is from Level~2 to Level~3 (54 percentage points), corresponding to the transition from problems requiring a single non-trivial theorem to those requiring deep synthesis of multiple advanced results.

The timeout rate provides an additional lens: at Level~3, nearly half the problems cause the model to fail without producing any response, compared to fewer than 5\% at Level~1.
This suggests that model confidence (or generation feasibility) correlates strongly with problem difficulty.

We note that difficulty calibration could also be influenced by the frequency of similar problems in training data.
Level~1 problems (e.g., ``give an example of a non-reduced scheme'') may appear frequently in online educational materials, while Level~3 problems are closer to research-level constructions with fewer training exemplars.

\subsection{Chapter-Level Analysis}
\label{sec:chapter-analysis}

Performance varies across chapters:

\begin{center}
\begin{tabular}{lcccc}
\toprule
\textbf{Chapter} & \textbf{$n$} & \textbf{Accuracy} & \textbf{Partial rate} & \textbf{Timeout rate} \\
\midrule
Cohomology & 15 & 0.93 & 0\% & 6.7\% \\
Basic Concepts & 47 & 0.79 & 4.3\% & 17.0\% \\
Computations & 21 & 0.76 & 0\% & 23.8\% \\
Moduli & 17 & 0.65 & 11.8\% & 23.5\% \\
\bottomrule
\end{tabular}
\end{center}

Several patterns emerge:
\begin{enumerate}[nosep]
  \item \textbf{Cohomology} has the highest accuracy (93\%), possibly reflecting the structured, algorithm-like nature of cohomological computations (spectral sequences, long exact sequences) that aligns well with sequential reasoning.
  \item \textbf{Basic Concepts} has moderate accuracy (79\%) despite being the most elementary topic, partly because this chapter contains many Level~2 problems requiring non-trivial constructions beyond textbook definitions.
  \item \textbf{Computations} has a 24\% timeout rate, the highest among all chapters, reflecting the difficulty of multi-step algebraic calculations (Chern classes, intersection products, Hodge numbers).
  \item \textbf{Moduli \& Deformation} has the lowest accuracy (65\%) and highest partial-credit rate (12\%), suggesting that the model struggles with the abstract constructions required in deformation theory (e.g., computing tangent spaces to moduli functors, constructing versal deformations).
\end{enumerate}

\subsection{Example vs.\ Counterexample Performance Gap}
\label{sec:gap-analysis}

The 15-percentage-point gap between example accuracy (80\%) and counterexample accuracy (65\%) is a key finding of the pilot evaluation.

\paragraph{Structural differences.}
Among partial-credit responses, 3 of 4 were on counterexample problems, suggesting that counterexample construction more often leads to partially correct responses where the right idea is present but execution is flawed.
The failure mode on these partial counterexample responses was predominantly F6 (incomplete verification) and F2 (property confusion): the model identified the relevant mathematical territory but could not precisely delineate the boundary between a theorem's hypotheses and its conclusion.

\paragraph{Timeout distribution.}
Among the 18 timeouts, 13 were on example problems and 5 on counterexample problems.
As a fraction of each type: 13/70 = 18.6\% of examples timed out vs.\ 5/23 = 21.7\% of counterexamples.
The slightly higher counterexample timeout rate, combined with the lower accuracy among completed responses, suggests that counterexamples are harder for the model both in terms of producing a response and in terms of producing a correct one.

\paragraph{Qualitative observation.}
In successful counterexample constructions, the model typically follows a clear pattern: (1)~state the target property, (2)~identify the critical hypothesis, (3)~construct an object where that specific hypothesis fails, and (4)~verify both that the object is valid and that the property indeed fails.
In unsuccessful attempts, the model often conflates the hypothesis and conclusion, or constructs an object that is loosely related to the topic without precisely targeting the boundary.

This pattern is consistent with the Lakatosian view~\citep{lakatos1976proofs} that counterexample construction requires ``concept stretching''---pushing a definition to its limits to find where it breaks.

\subsection{Limitations of Self-Evaluation}
\label{sec:self-eval-limitations}

A significant limitation of this pilot is that both the problem solver and the scorer are LLMs from the same model family (Claude Sonnet~4.6).
Among the 82 problems that received responses, 78 (95.1\%) were scored as correct.
This high concordance rate could reflect either genuine high accuracy or systematic bias in self-evaluation.

Several factors suggest caution:
\begin{itemize}[nosep]
  \item The scorer may share the solver's blindspots, consistently rating as correct a construction that both models find plausible but that a human expert would flag as flawed.
  \item The scorer was given the known solution, which may bias it toward confirming any response that follows a similar structure.
  \item With only 4 non-timeout failures among 82 responses, the failure mode taxonomy remains sparsely populated.
\end{itemize}

Future work will address this through human expert scoring of a stratified sample, cross-model evaluation (using different models as solver and scorer), and CAS-based automated verification where possible.
